\chapter{Amblyopia}

\section*{Definition}
Amblyopia, commonly called lazy eye, is reduced vision caused by abnormal visual development in childhood. The eye may look structurally normal, but the brain has not learned to process its visual input normally. It is usually unilateral but can affect both eyes. A fixed difference of two lines on a particular chart is not required for diagnosis; age-appropriate, best-corrected visual acuity and examination are interpreted together.

\section*{When to See a Doctor}

\begin{itemize}
 \item Child fails vision screening at school or pediatric well-visit
 \item Parents notice persistent head tilt, squinting, or eye misalignment
 \item Worsening visual acuity despite prior treatment
 \item New double vision, a new eye turn, or sudden vision loss at any age -- these are not typical of childhood amblyopia and require prompt assessment
\end{itemize}

\section*{How It Develops}
During early visual development, the brain favors one eye when images are persistently blurred, misaligned, or blocked. Connections that process the less-used eye are then underdeveloped. Visual development is most sensitive in early childhood, but some older children and adolescents can still improve with treatment.

\section*{Causes}
\begin{itemize}
 \item \textbf{Strabismic:} Esotropia or exotropia causing cortical suppression to avoid diplopia
 \item \textbf{Refractive:} High anisometropia (asymmetric refractive error) producing blurred retinal image in one eye
 \item \textbf{Deprivation:} Cataract, ptosis, or corneal opacity obstructing the visual axis during development
 \item \textbf{Combined:} Coexisting strabismic and refractive mechanisms
\end{itemize}

\section*{Related Symptoms}
\begin{itemize}
 \item Strabismus (esotropia, exotropia)
 \item Anisometropia
 \item Reduced depth perception or difficulty judging distances
 \item Squinting, closing one eye, or turning the head to see
 \item Poor visual performance in one eye, sometimes without obvious symptoms
\end{itemize}
\textbf{Important:} Early detection improves the chance of recovery, but treatment may still help some older children and adolescents. A new vision change in an adult should not be assumed to be amblyopia.

\section*{Medical Evaluation}
\begin{itemize}
 \item Comprehensive ophthalmic examination including cycloplegic refraction
 \item Best-corrected visual acuity testing (age-appropriate method: Lea symbols, HOTV, or Snellen)
 \item Cover-uncover and alternate cover test to detect strabismus
 \item Retinal examination to exclude structural causes of vision loss
 \item Additional testing, referral, or imaging when the examination suggests cataract, retinal disease, optic-nerve disease, or another structural cause; a neutral-density filter is not a routine confirmatory test
\end{itemize}

\section*{Treatment Options}
\begin{itemize}
 \item Full, age-appropriate optical correction of refractive error is usually the first treatment, followed by reassessment after an appropriate period of spectacle wear
 \item Occlusion therapy with a patch over the better-seeing eye for a prescribed number of hours, adjusted to severity and response
 \item Pharmacologic penalization with prescribed atropine in the better-seeing eye can be an alternative to patching for selected children; the dose and schedule must come from the eye specialist
 \item Bangerter filter foils for partial occlusion in mild amblyopia
 \item Digital or dichoptic therapies may be considered as adjuncts, but they do not replace spectacles, patching, or atropine when those are indicated
 \item Treat the cause of deprivation promptly, sometimes with surgery for a cataract, ptosis, or other obstruction. Strabismus surgery may be needed for alignment or binocular function but does not directly replace amblyopia treatment.
\end{itemize}

\section*{Self-Care and Home Management}
\begin{itemize}
 \item Compliance with prescribed occlusion therapy (eye patching) for the recommended daily duration
 \item Use atropine only as prescribed; do not change the eye, dose, or schedule, and keep the medicine away from other children
 \item Ensuring consistent spectacle wear for any refractive component
 \item Attend follow-up at the interval set by the ophthalmology or orthoptic team; visual acuity and treatment response must be monitored
\end{itemize}

\section*{References}
\begin{thebibliography}{99}
\bibitem{amblyopia:ref1} Holmes JM, Clarke MP. ``Amblyopia.'' \textit{Lancet}. 2006;367(9519):1343-1351.
\bibitem{amblyopia:ref2} Wallace DK, Repka MX, et al. ``Treatment of amblyopia in children 7 to 17 years old.'' \textit{Ophthalmology}. 2018;125(8):1171-1181.
\bibitem{amblyopia:ref3} American Academy of Ophthalmology. ``Preferred Practice Pattern: amblyopia.'' AAO, 2022.
\bibitem{amblyopia:ref4} Schachat AP, Wilkinson CP, et al. \textit{Ryan\'s Retina}, 7th ed. Elsevier, 2022.
\bibitem{amblyopia:ref5} UpToDate. ``Amblyopia in children: management.'' Wolters Kluwer, 2023.
\bibitem{amblyopia:ref6} American Academy of Ophthalmology. Amblyopia Preferred Practice Pattern. Updated 2022.
\end{thebibliography}
